% !PRINTED 292
\weberpage{292}
If $\beta \equiv 0$, then (5) always has the one root $z = \infty$; if
moreover $\delta \equiv \alpha$, and $\gamma$ not $\equiv 0$ $(\mathrm{mod}\ p)$, then there is only this
one; in this case, however,
\[
\alpha\delta \equiv 1, \qquad \alpha \equiv \delta, \qquad \left(\frac{\alpha+\delta}{2}\right)^2 \equiv 1 \quad (\mathrm{mod}\ p),
\]
and condition I. is satisfied. But if at the same time $\gamma \equiv 0$ $(\mathrm{mod}\ p)$,
then the substitution is the \gesperrt{identical one}.

If $\beta \equiv 0$, but $\alpha-\delta$ not $\equiv 0$ $(\mathrm{mod}\ p)$, then (5) has yet
a second root; since now $\alpha\delta \equiv 1$ $(\mathrm{mod}\ p)$, we have:
\[
\left(\frac{\alpha+\delta}{2}\right)^2 - 1 \equiv \left(\frac{\alpha-\delta}{2}\right)^2 \quad (\mathrm{mod}\ p),
\]
hence a quadratic residue, and condition II. is satisfied. We
summarise these theorems as follows:

\gesperrt{In case I. one element or all elements remain
unchanged, in case II. two elements remain unchanged
 and in case III. all elements are
changed.}

When we repeat one and the same substitution
\weq{7}{A = \begin{pmatrix}\alpha, & \beta \\ \gamma, & \delta\end{pmatrix}}
several times, the substitutions
\[
A, \ A^2, \ A^3 \dots,
\]
arise, in whose sequence the identical substitution $\begin{pmatrix}1, & 0 \\ 0, & 1\end{pmatrix}$
must once occur. If $n$ is the least positive number for which
\[
A^n = \begin{pmatrix}1, & 0 \\ 0, & 1\end{pmatrix}
\]
holds, then $n$ is called the degree of $A$ (Bd. II, \S\,2).

If we set, for an arbitrary $m$,
\[
A^m = \begin{pmatrix}\alpha_m, & \beta_m \\ \gamma_m, & \delta_m\end{pmatrix},
\]
then for the computation of the numbers $\alpha_m$, $\beta_m$, $\gamma_m$,
$\delta_m$ we obtain the following system of recurrent formulae:
\weq{8}{\begin{aligned}
\alpha_{m+1} &= \alpha\alpha_m+\beta\gamma_m, & \beta_{m+1} &= \alpha\beta_m+\beta\delta_m,\\
\gamma'_{m+1} &= \gamma\alpha_m+\delta\gamma_m, & \delta_{m+1} &= \gamma\beta_m+\delta\delta_m.
\end{aligned}} % [sic] printed with a prime on \gamma'_{m+1}, inconsistent with the other three terms

From these the numbers $\alpha_m$, $\beta_m$,
$\gamma_m$, $\delta_m$ can be computed with particular ease in case I.
